\section{Conclusão}

Concluímos que a arquitetura de recuperação por BM25 e trigramas seguida de reordenação supervisionada melhora a priorização de nomes de logradouros para o catálogo e as transformações avaliados no Pauliceia. Os trigramas ampliam a cobertura dos candidatos, e o Random Forest corrige as 36 falhas de primeira posição restantes no teste reservado. A comparação adicional indica ganho de 17 acertos sobre Token Set, o melhor comparador isolado nesse conjunto.

O empate com Damerau--Levenshtein no complemento mostra que o benefício adicional depende das características das consultas. A contribuição é uma implementação integrada e uma avaliação reproduzível dessa combinação, com distinção entre correspondência textual e comprovação histórica. A validação local confirma a entrega das sugestões e dos pontos cadastrados nas condições testadas. A extensão das conclusões para uso real exige consultas e correspondências históricas verificadas de forma independente.
